\chapter{How to read this manual}
\label{ch:howto}

This manual is written for scientists who want to prepare a protein for a
computer simulation, not for programmers. You do not need to know how to
program. You will need to type a few short commands, and this manual shows you
each one exactly.

It describes \textbf{ProPrep version 1.23.0}. Every step in the step-by-step
chapters was carried out on a real computer while this manual was being
written, and the screens shown are the screens ProPrep actually printed. Where a
feature could not be tried (because it needs software we did not have), the
text says so.

\section{How ProPrep talks to you}

ProPrep is a conversation. It shows you a short menu or asks a question, you
type an answer and press \key{Enter}, and it moves on. Most answers are a single
letter or a short code such as \key{1a} or \key{y}. When a question shows a
value in brackets, that value is the \term{default}: pressing \key{Enter}
without typing anything accepts it.

\begin{screen}
Run BLAST with these parameters? [y/n] (y): §ans«y»
\end{screen}

Here \prompt{[y/n]} tells you the possible answers, \prompt{(y)} is the
default, and the highlighted \screenans{y} is what was typed.

\section{What the different styles mean}

\begin{tabularx}{\linewidth}{@{}lL@{}}
\toprule
\textbf{Looks like} & \textbf{Means} \\
\midrule
\key{4e} & A key or menu code you type in ProPrep, then press \key{Enter}. \\
\typed{1UBQ} & Text you type yourself, such as a name or a number. \\
\prompt{Enter your choice} & A question or message exactly as ProPrep shows it. \\
\file{filtered_structure.pdb} & The name of a file or folder on your computer. \\
\menupath{4e}{Protonation State Analyzer} & A tool in ProPrep's main menu, with its code. \\
\pdb{1UBQ} & The four-character code of a structure in the Protein Data Bank. \\
\term{topology} & A technical word explained where it first appears and in the Glossary. \\
\bottomrule
\end{tabularx}

\medskip
A dark box shows what ProPrep prints on the screen; long screens are shortened,
and \texttt{...} marks where lines were left out:

\begin{screen}
══ 4. STRUCTURE PREPARATION ══
  4a. ○ PDB Filter
        Filter PDB structures by component
        ⚠ Needs a loaded structure
...
Enter your choice (1a): §ans«1a»
\end{screen}

A white box shows a command you type in your computer's \term{Terminal}
(the program where you type commands), \emph{outside} ProPrep:

\begin{terminal}
proprep --version
\end{terminal}

\section{The colored boxes}

\begin{tip}
A shortcut or a good habit that saves time.
\end{tip}

\begin{watchout}
A trap that caught us while testing ProPrep, and how to avoid it. Read these;
several of them can cost you an hour.
\end{watchout}

\begin{needsoftware}
This step needs a program that is not part of ProPrep, such as Gaussian or a
MODELLER license key. The box says what is needed and what happens without it.
\end{needsoftware}

\begin{behindscenes}
Optional background: what ProPrep is doing for you and why. You can skip these.
\end{behindscenes}

\section{How the step-by-step chapters are organized}

Each chapter in Part~IV starts with an \textsf{At a glance} box that tells you
what the tool does, when you need it, what must be done first, and which files
it produces. Then comes \textsf{Step by step}, the recommended route that worked
in our tests, followed by \textsf{Other choices} and \textsf{If something goes
wrong}.

\begin{glance}
\glanceitem{What it does}{One or two sentences.}
\glanceitem{You need it when}{The situations where this step matters.}
\glanceitem{Before you start}{Steps or software that must come first.}
\glanceitem{Where to find it}{\key{4e} in the main menu.}
\glanceitem{You get}{The files it writes.}
\end{glance}

Numbered steps look like this:

\begin{steps}
  \item At \prompt{Enter your choice}, type \key{1a} and press \key{Enter}.
  \item Type the PDB code, for example \typed{1UBQ}.
\end{steps}

\section{Where to start}

\begin{itemize}
  \item \textbf{New to ProPrep?} Read Chapters~\ref{ch:requirements}--\ref{ch:starting},
        then work through the tutorial in Chapter~\ref{ch:tutorial}.
  \item \textbf{Know what you want to do?} The table in Chapter~\ref{ch:bigpicture}
        lists the chapters you need for each kind of project.
  \item \textbf{Something went wrong?} Go to Chapter~\ref{ch:troubleshooting}.
\end{itemize}
